\chapter{A criação da Caixa de Conversão, 1906}

\section{O prelúdio, a defesa do café em 1905}

Quando os governos de São Paulo, Minas Gerais e Rio de Janeiro retomaram a discussão de uma ação conjunta em 1905, a crise do café já acumulava quase uma década. O problema deixou de ser apenas o baixo preço corrente e passou a incluir o tamanho da safra esperada, os estoques existentes e a incapacidade dos produtores de coordenar a oferta. Reuniões entre representantes estaduais e federais durante o segundo semestre procuraram obter apoio para um empréstimo que financiasse a retenção do produto \cite{holloway1978,torelli2004}.

A imprensa registrou uma combinação que permaneceria importante no ano seguinte: apoio à intervenção podia coexistir com desacordo sobre quem deveria executá-la e assumir os riscos. O \textit{Correio da Manhã} defendia ação diante da crise, mas, em novembro de 1905, preferia uma empresa privada apoiada pelos governos à operação mercantil direta dos estados. Cândido Rodrigues, em conferência na Sociedade Nacional de Agricultura, atacou o \textit{laissez-faire} e defendeu compras capazes de sustentar preços. No fim do ano, uma autorização orçamentária removeu obstáculos à garantia federal de empréstimos estaduais, sem fazer a União absorver todo o risco.

\fonteaconferir{Conferir na imagem o editorial de Gil Vidal, \textit{Correio da Manhã}, 21 de novembro de 1905, e a cobertura da conferência de Cândido Rodrigues em 1º de dezembro. O texto acima deriva da monografia, p. 29 a 31.}

Esse prelúdio mostra que a defesa do café não produzia automaticamente uma posição monetária. A conexão entre valorização do produto e fixação do câmbio ganhou forma quando os organizadores calcularam quanto capital externo seria necessário e perceberam que o valor do empréstimo, expresso em libras e convertido em mil-réis, dependia da taxa. A apreciação também poderia retirar, em moeda nacional, parte do ganho obtido com a elevação do preço externo.

\section{Taubaté, a união de dois problemas}

O Convênio de Taubaté foi assinado na madrugada de 26 de fevereiro de 1906 pelos presidentes dos três principais estados cafeeiros. Previa compras e retenção de estoques, sobretaxa por saca, restrições a novas plantações e um empréstimo externo. A criação de uma caixa que fixasse o valor da moeda aparecia como condição do esquema. Café e câmbio se tornaram, naquele momento, componentes de uma mesma proposta política, embora exigissem competências e instrumentos jurídicos diferentes \cite{holloway1978,torelli2004}.

A autoria da cláusula cambial foi ela própria objeto de disputa. O \textit{Correio da Manhã} atribuiu a indicação a Nilo Peçanha. Edmundo Bittencourt sustentou depois que a exigência teria servido para dificultar a valorização, pois Nilo desconfiava dos riscos e da vantagem desproporcional de São Paulo. Augusto Ramos ofereceu outra explicação: os delegados teriam chegado à fixação ao estimar o empréstimo necessário sob taxas que variavam entre 12 e 18 pence. Em debates posteriores, Toledo Piza e Pinheiro Machado também apareceram como responsáveis pela entrada da ideia na agenda.

\fonteaconferir{Conferir a notícia do \textit{Correio da Manhã} de 27 de fevereiro, o artigo de Edmundo Bittencourt de 13 de junho e a carta de Augusto Ramos publicada pela \textit{Gazeta de Notícias} em 6 de setembro. A divergência não deve ser resolvida sem os documentos do Convênio e a carta de Pinheiro Machado a Jorge Tibiriçá.}

A disputa de paternidade importa porque atribuir a Caixa a Nilo, aos técnicos paulistas, à viagem argentina ou à articulação de Pinheiro Machado modificava seu sentido político. A instituição podia aparecer como condição técnica do empréstimo, salvaguarda de um estado relutante, projeto de proteção da produção ou instrumento de uma coalizão parlamentar. Em vez de escolher uma origem única, a dissertação tratará essas versões como apropriações concorrentes.

Nos primeiros dias de março, a \textit{Gazeta de Notícias} apoiou o objetivo de estabilizar a moeda, mas manifestou receio diante do desenho do Convênio. O jornal temia a circulação simultânea de papel inconversível e notas conversíveis emitidas para financiar a valorização. Poucos dias depois, ao comentar carta de Custódio Coelho, formulou a controvérsia entre o par de 27 e a proposta de 12 como espaço aberto a combinações. O \textit{Correio Paulistano}, por outro lado, respondeu às acusações de quebra do padrão dizendo que o papel já existente não seria compulsoriamente convertido a uma taxa reduzida. O novo aparelho abriria entrada ao capital e protegeria produção e comércio contra oscilações.

\fonteaconferir{Conferir os editoriais da \textit{Gazeta de Notícias} de 5, 7 e 9 de março e a coluna do \textit{Correio Paulistano} de 15 de março. As passagens são essenciais para distinguir apoio ao objetivo, crítica ao desenho e posição sobre a paridade.}

\section{Rodrigues Alves e a separação dos projetos}

Jorge Tibiriçá pediu a Rodrigues Alves que convocasse sessão extraordinária do Congresso para aprovar o Convênio. O presidente recusou. A política monetária pertencia à União, e a fixação em taxa inferior ao padrão legal colidia, em sua avaliação e na de Leopoldo de Bulhões, com o caminho adotado desde 1898. A mensagem presidencial de maio enviou o acordo ao Congresso, mas criticou a ideia de que a prosperidade da lavoura dependesse de câmbio baixo.

O \textit{Correio da Manhã}, apesar de oposicionista, elogiou a recusa presidencial e associou o Convênio a aventura, quebra do padrão e compromissos com credores internos e externos. A passagem mostra como a posição diante de um ato podia atravessar alinhamentos governistas mais gerais. Fora do corpus principal atual, \textit{O Estado de S. Paulo} criticou a recusa e acusou Rodrigues Alves de condenar o Convênio sob aparência de neutralidade. O contraste preservado na monografia será usado apenas como evidência acessória, com estatuto explícito.

\fonteaconferir{Conferir o editorial do \textit{Correio da Manhã} de 11 de março de 1906. A menção a \textit{O Estado de S. Paulo}, 16 de março, vem do corpus da monografia e não integra o censo principal desta dissertação.}

Em junho, os articuladores separaram a valorização do café da fixação cambial. A primeira podia prosseguir por acordos e autorizações relacionados aos estados; a segunda exigia lei federal e reorganização monetária. A separação reduziu a dependência imediata entre os projetos, mas não apagou sua origem comum. A valorização efetivamente implementada acabou mais concentrada em São Paulo e apoiada por negociantes e bancos estrangeiros do que o Convênio original sugeria \cite{holloway1978}.

David Campista, deputado mineiro e futuro ministro da Fazenda de Afonso Pena, tornou-se relator e principal formulador do projeto da Caixa. A bancada mineira, representantes paulistas e o Bloco articulado por Pinheiro Machado sustentaram a aprovação. A oposição, contudo, não formava um campo único. Essa composição é o primeiro indício de que um eixo simples entre apoio e rejeição perde diferenças doutrinárias relevantes.

\section{Três correntes e um vocabulário comum}

Uma primeira corrente defendia a Caixa proposta por Campista como barreira à valorização excessiva e à especulação. Entre seus representantes estavam o próprio relator, Pinheiro Machado, Adolpho Gordo, Altino Arantes e Barros Franco Júnior. O argumento ligava estabilidade à possibilidade de cálculo, ao ingresso de capitais e à proteção da produção nacional. O café aparecia tanto como interesse setorial quanto como principal fonte de receita externa, o que permitia apresentar sua defesa como interesse nacional \cite{torelli2007}.

Uma segunda corrente preservava o horizonte de Murtinho e Bulhões. Barbosa Lima, Affonso Costa e Paula Ramos criticavam a fixação abaixo do par legal, o regionalismo e o risco de reabrir inflação. A política deveria conservar fundos de resgate e garantia, limitar emissões e permitir que a valorização prosseguisse na direção de 27 pence. Parte do vocabulário desse campo associava a mudança do padrão a quebra contratual, perda de probidade e bancarrota.

Uma terceira corrente, representada por Alcindo Guanabara e Serzedello Corrêa, também se opunha ao projeto Campista, mas por razão diferente. Considerava irreal o retorno a 27 e propunha reconhecer um novo padrão, em torno de 12 pence, acompanhado de conversibilidade mais ampla e saneamento da circulação. Para esse grupo, a Caixa de 15 podia ser um organismo incompleto, incapaz de impedir movimentos nos dois sentidos e sujeito a inflação. Sua defesa de uma taxa menor não significava abandono do metalismo. Procurava instituir moeda conversível sobre uma paridade julgada compatível com a economia.

\fonteaconferir{Conferir a cobertura do discurso de Serzedello Corrêa no \textit{Correio da Manhã} de 31 de agosto e das emendas de Alcindo Guanabara nos anais e nos jornais. Conferir também a cobertura de Paula Ramos em 1º de setembro.}

Os três campos compartilhavam um vocabulário que ajuda a localizar o horizonte do debate: fixação, par legal, curso forçado, curso legal, circulação metálica, numerário, resgate, garantia, especulação e estabilidade. Divergiam sobre como esses termos se relacionavam. Para uns, aceitar 15 quebrava o padrão; para outros, o papel inconversível já impedia que o par de 27 tivesse realidade econômica; para os defensores do projeto, as notas novas formavam contrato distinto e não anulavam a obrigação futura relativa ao papel antigo.

\section{A taxa de fixação: 12, 15 ou 27 pence}

A decisão institucional foi fixar o valor das notas da Caixa em 15 pence por mil-réis. A taxa aproximava-se do câmbio corrente e funcionava como compromisso. Não adotava a proposta cafeeira mais baixa, não decretava a conversão geral a 12 e interrompia, na prática, a valorização que poderia levar ao par de 27. Ao mesmo tempo, a lei preservou pagamentos contratados em ouro segundo o padrão legal e permitiu que o Congresso elevasse a taxa no futuro \cite{brasil1906caixa}.

O \textit{Correio Paulistano} publicou em 3 de julho a justificação técnica do projeto e, em 11 de outubro, sua redação aprovada. Como esses textos eram documentos reproduzidos, seu conteúdo prova circulação e destaque, não voz editorial automática. A apropriação do jornal aparece com maior segurança na coluna \textit{Política Republicana}, de 29 de agosto, que ligou a Caixa, a valorização e a ação do governo paulista a uma obra patriótica. A publicação do discurso de Adolpho Gordo acrescentou uma defesa segundo a qual a taxa poderia ser alterada depois, afastando a acusação de quebra definitiva.

\fonteaconferir{Conferir visualmente a justificação de 3 de julho, a coluna de 29 de agosto, o discurso publicado em 9 de setembro e o texto final de 11 de outubro no \textit{Correio Paulistano}.}

O \textit{Correio da Manhã} foi o adversário mais duro entre os jornais centrais já lidos para 1906. Elogiou oradores de correntes oposicionistas diferentes, atacou Campista por mudanças no projeto e utilizou o relatório de Custódio Coelho para acusar incoerência entre a posição do técnico em março e a de novembro. O apoio a Paula Ramos aproximava o jornal da continuidade deflacionária, enquanto o elogio a Serzedello mostra que a hostilidade ao projeto podia hospedar uma alternativa distinta do retorno a 27.

\fonteaconferir{Conferir as crônicas do \textit{Correio da Manhã} de 25 de setembro e 10 de novembro, inclusive título, seção, autoria e relação entre comentário do jornal e discurso reproduzido.}

A \textit{Gazeta de Notícias} percorreu um caminho que também exige decomposição. Em março apoiou a estabilização, criticou a possibilidade de duas moedas e procurou uma solução entre 12 e 27. Em novembro, depois da posse de Afonso Pena, declarou apoio decidido ao projeto aprovado e felicitou Congresso e ministro. A sequência pode representar mudança do desenho avaliado, mudança política ou combinação das duas. Só a leitura completa das peças intermediárias permitirá escolher.

\fonteaconferir{Conferir o editorial da \textit{Gazeta de Notícias} de 27 de novembro e reconstruir a continuidade entre a crítica de março e o apoio final.}

No estado atual da leitura, \textit{O Paiz} ainda precisa de uma reconstrução sistemática para 1906. A ausência de uma conclusão antecipada é preferível a transferir ao jornal o governismo geral ou as posições de vozes que ele reproduziu.

\lacuna{Completar a janela de fevereiro a dezembro de 1906 em \textit{O Paiz}, distinguindo editoriais, relatos do Congresso, seções livres e documentos transcritos.}

\section{O limite da emissão}

O projeto vinculava cada nota emitida a ouro recebido. Ainda assim, o limite total importava porque determinava quanto ouro poderia ser absorvido à taxa fixada antes de nova decisão. Um teto estreito permitiria que, uma vez atingido, o câmbio retomasse a valorização. Um teto amplo prolongaria a fixidez e acomodaria maior expansão da circulação.

A lei estabeleceu 320 mil contos, correspondentes a 20 milhões de libras. Ao alcançar esse montante, as emissões cessariam e apenas nova lei poderia elevar a taxa. A cifra de 320 mil contos já aparecia na justificação publicada em julho, mas o relatório de pesquisa registra ali correspondência com 15 milhões de libras. A divergência pode refletir uma versão anterior, erro de transcrição ou relação cambial distinta e precisa ser conferida antes de ser interpretada.

\fonteaconferir{Comparar na imagem o texto do \textit{Correio Paulistano} de 3 de julho com a redação de 11 de outubro e com o art. 3º da lei. Não explicar a diferença entre 15 e 20 milhões de libras antes dessa conferência.}

O vocabulário do debate ligava limite a inflação, elasticidade e continuidade da valorização. O argumento favorável afirmava que a saída de ouro recolheria as notas na mesma proporção e impediria emissão sem lastro. O argumento contrário questionava não apenas a proporção um para um, mas o efeito de grande entrada de ouro sobre o volume total de papel e sobre o caminho até o par.

\section{Os fundos de resgate e garantia}

Os fundos instituídos em 1899 corporificavam a política anterior. O de resgate destinava recursos à retirada de papel; o de garantia acumulava ouro e divisas para preparar a conversibilidade e apoiar operações cambiais. Preservá-los intactos significava conservar o aparelho da valorização gradual. Transferi-los à nova Caixa significava incorporar recursos da ortodoxia de 1899 ao regime de taxa nova.

A lei transferiu ambos à Caixa. Os saldos do fundo de resgate continuariam destinados ao fim original. O fundo de garantia também poderia resgatar papel, permutado por notas conversíveis, e até três milhões de libras poderiam apoiar operações do governo para manter a taxa \cite{brasil1906caixa}. A solução não simplesmente extinguiu a herança de 1899. Reorientou seus instrumentos para a estabilidade a 15.

Essa decisão ajuda a explicar por que a Caixa não cabe numa oposição entre mecanismo heterodoxo e ortodoxia anterior. A nova instituição absorveu fundos criados para o retorno ao regime metálico, emitiu apenas contra ouro e preservou formalmente o padrão de 27 para obrigações especificadas. O conflito incidia sobre o uso desses recursos e sobre a trajetória futura do câmbio.

\section{Quem poderia alterar a taxa}

Outra controvérsia dizia respeito à autoridade. Delegar ao Executivo a mudança cambial permitiria resposta mais rápida às condições do mercado, mas entregaria ao governo poder sobre contratos, emissão e distribuição de ganhos. Reservar a decisão ao Congresso tornava a mudança mais lenta e pública.

O texto final adotou uma divisão. A Caixa ficou sob superintendência imediata do ministro da Fazenda, e o Executivo recebeu poder para regulamentar a administração e realizar operações cambiais. A taxa, porém, só poderia ser elevada por lei depois de atingido o limite dos depósitos. A decisão de 1906 preservou, portanto, a competência legislativa sobre a paridade, problema que voltaria com força em 1910.

\fonteaconferir{Localizar nos \textit{Documentos Parlamentares} os argumentos explícitos a favor da competência privativa do Congresso e da delegação ao Executivo.}

\section{A agência em Londres e o controle do lastro}

O projeto autorizou a criação de uma agência da Caixa em Londres. A solução poderia evitar o transporte físico de todo o ouro e facilitar depósitos, saques e operações no principal centro financeiro ligado ao Brasil. Também levantava dúvidas sobre a localização do lastro, a autoridade dos agentes financeiros e a capacidade de fiscalização nacional.

A lei permitiu que a agência emitisse notas conversíveis à vista e a colocou sob superintendência do ministro. Autorizou ainda o governo a comprar e vender letras para manter a taxa, inclusive por seção especial do Tesouro, mas excluiu a própria Caixa dessas operações discricionárias \cite{brasil1906caixa}. A separação protegia formalmente o ouro recebido contra uso diverso do resgate, enquanto criava uma via paralela de intervenção cambial.

\fonteaconferir{O contraditório sobre a agência de Londres está pouco documentado no material já lido. Reunir discursos, emendas de setembro e comentários dos jornais antes de atribuir polos definidos.}

\section{A lei e o compromisso de 1906}

A Lei n. 1.575, de 6 de dezembro, instituiu notas com curso legal, resgatáveis à vista e garantidas por depósito equivalente. O ouro não poderia ser usado para outra finalidade, e as notas resgatadas seriam inutilizadas. A Caixa publicaria mensalmente o estado de depósitos e emissões. Pagamentos já devidos em ouro continuariam referidos a 27 pence, embora pudessem ser realizados em notas da Caixa pelo valor metálico representado. O desenho criou uma circulação conversível dentro de um sistema no qual o papel do Tesouro continuava inconversível \cite{brasil1906caixa,oliveirasilva2001}.

O compromisso reuniu elementos que os polos apresentavam como incompatíveis. Fixou 15, mas não revogou formalmente o padrão de 27. Permitiu emissão elástica à entrada de ouro, mas exigiu lastro integral e teto. Incorporou os fundos de 1899, mas modificou sua função. Entregou administração e intervenção ao Executivo, mas reservou a elevação da taxa ao Congresso. Autorizou agência externa, mas criminalizou o desvio dos depósitos.

\interpretacaoprovisoria{A aprovação não encerrou a controvérsia, apenas transformou suas perguntas. Depois de dezembro, já não se discutia uma Caixa abstrata. Tornou-se possível observar se a taxa de 15 se sustentaria, se o limite seria alcançado, como o Banco do Brasil interviria, que uso seria feito dos fundos e qual parte do compromisso sobreviveria a uma mudança da conjuntura externa.}

